% Book 3, Year 1, chapter 30 -- project: a search engine for a folder of texts, in C.
% Provenance (never printed): sources/book3/BRIEFS.md, chapter 30.
% ledger: b1:tfidf, b1:zipf
\chapter{Project: A Search Engine for a Folder of Texts}\label{ch:b1:project-text-search-engine}

A folder holds four hundred texts, 1.9 megabytes in all, and someone wants to know which of them mention
\emph{quartzite}. Reading every file answers the question in a few milliseconds. Now the same person asks a second
question, then a hundred, then a thousand, and each time the program reads the whole folder again, although the texts
have not changed. A search engine reads them once, writes down where every word occurs, and answers each question from
that record. This chapter builds one in C, in six milestones, from a specification to a measurement. Nothing in it is
new: it puts together the files of \cref{ch:b1:c-strings-structs-files}, the modules and tests of
\cref{ch:b1:c-modules-make-tests}, the history of \cref{ch:b1:version-control}, the \omterm{def:b1:dynamic-arrays:dynamic-array}{dynamic arrays} of
\cref{ch:b1:dynamic-arrays}, the \omterm{def:b1:hash-tables:hash-table}{hash tables} of \cref{ch:b1:hash-tables} and the heaps of \cref{ch:b1:heaps}. What is
new is making them work together, and showing that they do.

\begin{recall}
A file is read through a stream, in binary mode when every byte matters, and a robust reader checks every call
(\cref{sec:b1:c-strings-structs-files:robust}). A module is a header (its interface) and a \omterm{def:b1:compilation-pipeline:source-file}{source file}; a \omterm{def:b1:c-modules-make-tests:makefile}{makefile}
rebuilds what is out of date (\cref{def:b1:c-modules-make-tests:makefile}). A \omterm{def:b1:dynamic-arrays:dynamic-array}{dynamic array} doubles its capacity when
it is full, at amortised constant cost per push (\cref{def:b1:dynamic-arrays:amortised}). A \omterm{def:b1:hash-tables:hash-table}{hash table} with \omterm{def:b1:hash-tables:collisions}{separate
chaining} finds a key in expected constant time while its \omterm{def:b1:hash-tables:load-factor}{load factor} stays bounded
(\cref{def:b1:hash-tables:load-factor}). The $k$ largest elements of a stream are kept in a min-heap of size $k$
(\cref{lst:b1:heaps:topk}).
\end{recall}

\section{Specification and corpus}\label{sec:b1:project-text-search-engine:spec}

The program is called \code{search}. Run as \code{search DIR WORD...}, it reads every file of the folder \code{DIR}
whose name ends in \code{.txt}, and prints how many of them contain \emph{every} word of the query, then the ten best of
them, best first, each with its score; a word is matched whatever its case. Run with a folder and no word, it prints the
size of what it built. Its \omterm{def:b1:the-shell:exit-status}{exit status} says what happened: 0 when some document matches, 1 when none does, 2 when the
command is misused or the folder cannot be read, so that a script can tell ``nothing found'' from ``something broke''.
It is built by a \omterm{def:b1:c-modules-make-tests:makefile}{makefile}, tested by \omterm{def:b1:c-modules-make-tests:unit-test}{unit tests} and by a comparison with a simpler program, kept in a Git \omterm{def:b1:version-control:repository}{repository},
and measured against that simpler program, which reads every file for every query.

\omimg{book3/ai/b1-project-text-search-engine-desk.jpg}{0.5\linewidth}{Four hundred texts are a desk of folders:
searching it once by reading everything is fine; searching it a thousand times is not.}

\paragraph{The corpus.} The texts are generated, not downloaded, by a script seeded so that every run writes the same
400 files. Each document has between 200 and 2000 words drawn from an invented vocabulary of 5000 words made of
syllables. The word of rank $r$ is drawn with probability proportional to $1/r$, so a handful of short words fill
most of every text and most words are rare, as in real prose. Three words that the vocabulary cannot produce (its
syllables have no \emph{q}, \emph{x}, \emph{y} or \emph{z}) are planted in known documents, so that tests know the
answer in advance: \emph{quixotic} in two documents, \emph{xylograph} in one and \emph{quartzite} in four. One
document in fifty starts with \emph{Café}, whose last letter is not ASCII, and one in forty has no final newline.

\omcode{code/bachelor-1/30-project-text-search-engine/listings/make_corpus.py}{22}{39}{One document: words drawn with
weights $1/r$, the planted words, the traps, and lines of ten words}\label{lst:b1:project-text-search-engine:corpus}

\omoutput{code/bachelor-1/30-project-text-search-engine/outputs/make_corpus.txt}{The corpus: the word of rank $r$
occurs about $47\,500 / r$ times, so rank times count is nearly constant.}\label{out:b1:project-text-search-engine:corpus}

The product in the last column is the law that George Kingsley Zipf described in 1949 for the words of real texts:
in a large text, the frequency of a word is roughly inversely proportional to its rank. The script imposes it here on
purpose; it matters for the engine, because it says that a few words occur in every document and most occur in very
few.

\paragraph{The first version.} Before any index, the specification is met by the simplest program that could meet it,
\code{naive}: for each query, read every file, count each query word in it, and score the files that contain them all.
It is slow, but it is short enough to be checked by reading, and it will serve as the \emph{oracle} of the tests: the
fast program must print, for every query, exactly what the slow one prints.

\omcode{code/bachelor-1/30-project-text-search-engine/listings/scan.c}{11}{42}{The naive scan: every file, every word,
compared with every term}\label{lst:b1:project-text-search-engine:scan}

Both programs share the code that cuts a text into words, the score and the printing, so that a difference between
them can only come from the index. The whole design fits in one picture.

\begin{omfigure}
\begin{tikzpicture}[>={Stealth[length=1.8mm]}, font=\small,
  m/.style={draw, rounded corners=2pt, fill=omCode!8, minimum height=8mm, align=center, inner sep=3pt},
  d/.style={draw, dashed, minimum height=8mm, align=center, inner sep=3pt}]
\node[d] (dir) at (0,0) {folder\\\code{*.txt}};
\node[m] (fold) at (2.4,0) {\code{folder}\\list, sort};
\node[m] (tok) at (5.0,0) {\code{tokenize}\\words};
\node[m] (ix) at (7.7,0) {\code{index}\\hash table};
\node[d] (q) at (0,-2.0) {query\\words};
\node[m] (qr) at (5.0,-2.0) {\code{query}\\AND, tf--idf};
\node[m] (tk) at (7.7,-2.0) {\code{topk}\\heap of 10};
\node[m] (out) at (10.2,-2.0) {\code{main}\\print};
\draw[->] (dir) -- (fold);
\draw[->] (fold) -- (tok);
\draw[->] (tok) -- node[above, font=\scriptsize] {words} (ix);
\draw[->] (q) -- node[above, font=\scriptsize] {\code{split_words}} (qr);
\draw[->] (ix) -- node[left, font=\scriptsize] {posting lists} (qr.north east);
\draw[->] (qr) -- node[above, font=\scriptsize] {hits} (tk);
\draw[->] (tk) -- (out);
\node[m, fill=omThm!10] (sc) at (2.4,-3.6) {\code{scan}\\(the oracle)};
\draw[->, omThm] (fold) to[bend right=30] (sc.north west);
\draw[->, omThm] (sc.east) -| (tk.south);
\end{tikzpicture}
\par\smallskip{\small The modules of the search engine and the data that flows between them. The texts are read once,
to build the index; each query then touches only the \omterm{def:b1:project-text-search-engine:inverted-index}{posting lists} of its words. The naive scan reuses the \omterm{def:b1:project-text-search-engine:tokenizer}{tokenizer}
and the heap but reads every file for every query.}
\end{omfigure}

\section{Splitting files into words}\label{sec:b1:project-text-search-engine:tokenize}

\begin{definition}[Tokenizer]\label{def:b1:project-text-search-engine:tokenizer}
A \emph{tokenizer}\index{tokenizer} is a function that cuts a text into a sequence of \emph{tokens}, the units that the
rest of a program works with, by a rule that says which characters belong to a token and which separate tokens.
\end{definition}

Here a token is a word, and the rule is short: a word is a maximal run of ASCII letters, lowercased; every other byte
separates words. The rule works on bytes, so it never has to decode the text. Its cost is a documented choice about
the other alphabets. In UTF-8 the letter \emph{é} is the two bytes \code{0xC3 0xA9}, neither of which is an ASCII
letter, so \emph{Café} becomes the word \emph{caf}. That is wrong for a reader of French, and a real engine decodes the
text and folds accents. It is consistent, though: the query \emph{Café} is cut by the same function into the same
\emph{caf}, so it still finds its documents. A word longer than 63 letters is cut to its first 63, so that it fits in a
fixed buffer without overflowing it.

\omcode{code/bachelor-1/30-project-text-search-engine/listings/tokenize.c}{8}{24}{The \omterm{def:b1:project-text-search-engine:tokenizer}{tokenizer}: skip separators, then
copy letters, lowercased, until the next separator}\label{lst:b1:project-text-search-engine:next-word}

The bit trick on line 18 lowercases an ASCII letter: \emph{A} is \code{0x41} and \emph{a} is \code{0x61}, so setting bit
5 maps every capital to its small letter and leaves small letters alone. Note what the function does at the end of the
text: the second loop stops at \code{len} as well as at a separator, so the last word of a file without a final newline
is kept (\cref{exo:b1:project-text-search-engine:5} shows a \omterm{def:b1:project-text-search-engine:tokenizer}{tokenizer} that loses it).

The file is read whole into memory, in binary mode, by doubling a buffer until \code{fread} returns nothing more, and
every failure gives back \code{NULL} with nothing left allocated.

\omcode{code/bachelor-1/30-project-text-search-engine/listings/tokenize.c}{26}{48}{Reading a whole file into a growing
buffer}\label{lst:b1:project-text-search-engine:read-file}

The \omterm{def:b1:project-text-search-engine:tokenizer}{tokenizer} gets its own \omterm{def:b1:c-modules-make-tests:unit-test}{unit tests}, on inputs chosen for their edges: an apostrophe, two dashes, capitals, a
non-ASCII letter, a digit inside a word, a word too long for the buffer, and a text whose last byte is a letter.

\omcode{code/bachelor-1/30-project-text-search-engine/listings/units_test.c}{14}{29}{\omterm{def:b1:c-modules-make-tests:unit-test}{Unit tests} of the \omterm{def:b1:project-text-search-engine:tokenizer}{tokenizer} on
tricky inputs}\label{lst:b1:project-text-search-engine:tok-test}

\section{The inverted index}\label{sec:b1:project-text-search-engine:index}

\begin{definition}[Inverted index, posting list]\label{def:b1:project-text-search-engine:inverted-index}
Let a collection of documents be numbered $0, 1, \dots, N - 1$. The \emph{posting list}\index{posting list} of a word
is the list of the documents that contain it, in increasing order, each with the number of times the word occurs in
it; each element is a \emph{posting}. An \emph{inverted index}\index{inverted index} of the collection is a \omterm{def:b1:hash-tables:dictionary}{dictionary}
that maps each word occurring in it to its posting list.
\end{definition}

A document is a list of words; the index turns that around and gives, for each word, its documents, hence the name.
It is built in one pass. The documents are read in order of their names, so the numbers are the same on every run, and
each word of document $d$ either adds a posting $(d, 1)$ at the end of its list or, if the last posting is already
$(d, c)$, makes it $(d, c + 1)$. Since $d$ only grows, every list stays sorted without being sorted. The \omterm{def:b1:hash-tables:dictionary}{dictionary} is
a \omterm{def:b1:hash-tables:hash-table}{hash table} with \omterm{def:b1:hash-tables:collisions}{separate chaining}, hashed by FNV-1a, whose number of \omterm{def:b1:hash-tables:hash-table}{buckets} doubles when there are more words than
\omterm{def:b1:hash-tables:hash-table}{buckets}; each \omterm{def:b1:project-text-search-engine:inverted-index}{posting list} is a \omterm{def:b1:dynamic-arrays:dynamic-array}{dynamic array} of postings.

\omcode{code/bachelor-1/30-project-text-search-engine/listings/index.h}{6}{28}{The types of the index: a posting, a
\omterm{def:b1:project-text-search-engine:inverted-index}{posting list}, a \omterm{def:b1:hash-tables:hash-table}{bucket} entry, the index}\label{lst:b1:project-text-search-engine:index-h}

\omcode{code/bachelor-1/30-project-text-search-engine/listings/index.c}{41}{80}{Finding a word's entry, adding an
occurrence, and indexing a text}\label{lst:b1:project-text-search-engine:index}

\begin{omfigure}
\begin{tikzpicture}[font=\small, >={Stealth[length=1.6mm]},
  b/.style={draw, minimum width=8mm, minimum height=6mm},
  e/.style={draw, rounded corners=1.5pt, fill=omCode!8, minimum height=6mm, inner sep=2pt},
  p/.style={draw, minimum width=10mm, minimum height=6mm, font=\scriptsize, fill=omThm!8}]
\foreach \i in {0,...,4} {\node[b] (k\i) at (0,-\i*0.8) {}; \node[font=\scriptsize, anchor=east] at (-0.5,-\i*0.8) {\i};}
\node[font=\scriptsize, anchor=east] at (-0.5,-4.4) {$\vdots$};
\node[e] (e1) at (2.2,-0.8) {\emph{cat}};
\node[e] (e2) at (4.0,-0.8) {\emph{hat}};
\node[e] (e3) at (2.2,-2.4) {\emph{the}};
\node[e] (e4) at (2.2,-3.2) {\emph{a}};
\draw[->] (k1.center) -- (e1);
\draw[->] (e1) -- (e2);
\draw[->] (k3.center) -- (e3);
\draw[->] (k4.center) -- (e4);
\foreach \t/\x in {0\,{,}\,1/0, 1\,{,}\,1/1, 2\,{,}\,1/2} {\node[p] at (6.8+\x*1.05,0.1) {(\t)};}
\node[font=\scriptsize, anchor=east] (lc) at (6.2,0.1) {list of \emph{cat}};
\draw[->, omThm] (e1.north) to[bend left=25] (lc.west);
\foreach \t/\x in {0\,{,}\,2/0, 1\,{,}\,2/1} {\node[p] at (6.8+\x*1.05,-2.4) {(\t)};}
\draw[->, omThm] (e3.east) -- (6.25,-2.4);
\foreach \t/\x in {2\,{,}\,3/0} {\node[p] at (6.8+\x*1.05,-3.2) {(\t)};}
\draw[->, omThm] (e4.east) -- (6.25,-3.2);
\node[font=\scriptsize, align=left, anchor=west] at (6.25,-1.2) {postings (document, count),\\sorted by document};
\end{tikzpicture}
\par\smallskip{\small Part of the \omterm{def:b1:project-text-search-engine:inverted-index}{inverted index} of three small documents (those of
\cref{exo:b1:project-text-search-engine:2}). The \omterm{def:b1:pointers-and-arrays:array}{array} of \omterm{def:b1:hash-tables:hash-table}{buckets} holds \omterm{def:b1:pointers-and-arrays:pointer}{pointers} to chains of entries; \emph{cat} and
\emph{hat} share a \omterm{def:b1:hash-tables:hash-table}{bucket} (the \omterm{def:b1:hash-tables:hash-table}{bucket} numbers are made up). Each entry owns its word and its \omterm{def:b1:project-text-search-engine:inverted-index}{posting list}, a \omterm{def:b1:dynamic-arrays:dynamic-array}{dynamic array} of (document,
count) pairs.}
\end{omfigure}

How big is it? The function \code{index_memory} adds up what the index allocated: the \omterm{def:b1:pointers-and-arrays:array}{array} of \omterm{def:b1:hash-tables:hash-table}{buckets}, the entries,
the copies of the words, the posting \omterm{def:b1:pointers-and-arrays:array}{arrays} at their capacity (not their size, since the slack is allocated too) and the
file names. On the corpus, the engine prints its statistics when it is given no word:

\omoutput[24-25]{code/bachelor-1/30-project-text-search-engine/outputs/build.txt}{The index of the corpus: 5004 distinct
words, 193\,455 postings, 2.46 million bytes.}\label{out:b1:project-text-search-engine:stats}

The corpus has 434\,742 words, and the index has 193\,455 postings: each posting stands for all the occurrences of a word
in a document, which is why there are fewer than half as many. The index is a little bigger than the texts, nine tenths of it
posting \omterm{def:b1:pointers-and-arrays:array}{arrays}, almost a third of whose capacity is the slack that doubling leaves (\cref{exo:b1:project-text-search-engine:11}).
The 5004 words are the vocabulary of 5000, every one of which occurs somewhere, the three planted words, and the
fragment \emph{caf}.

\section{Answering queries}\label{sec:b1:project-text-search-engine:queries}

A query of $k$ words asks for the documents that contain all of them: the intersection of their \omterm{def:b1:project-text-search-engine:inverted-index}{posting lists}. Two
sorted lists are intersected by merging, as in \omterm{def:b1:sorting:merge-sort}{merge sort}: walk through both with one index each, advance the one that
points at the smaller document, and keep a document when both point at it. Every step advances at least one index, so
the merge of lists of lengths $\ell$ and $m$ costs at most $\ell + m$ steps. For $k$ words, the lists are taken
shortest first. The running result is then never longer than the shortest list, and it shrinks as the merges go on.

\begin{omfigure}
\begin{tikzpicture}[font=\small, c/.style={draw, minimum width=7mm, minimum height=6mm},
  k/.style={c, fill=omThm!25}, >={Stealth[length=1.5mm]}]
\node[anchor=east] at (-0.45,0) {word 1};
\node[anchor=east] at (-0.45,-1.3) {word 2};
\node[anchor=east] at (-0.45,-2.6) {both};
\foreach \v/\x/\s in {3/0/k, 8/1/c, 12/2/k, 19/3/k, 27/4/c} {\node[\s] (a\x) at (\x*0.85,0) {\v};}
\foreach \v/\x/\s in {1/0/c, 3/1/k, 4/2/c, 12/3/k, 15/4/c, 19/5/k, 30/6/c} {\node[\s] (b\x) at (\x*0.85,-1.3) {\v};}
\foreach \v/\x in {3/0, 12/1, 19/2} {\node[k] (r\x) at (\x*0.85,-2.6) {\v};}
\draw[<->, omCode] (a0.south) -- (b1.north);
\draw[<->, omCode] (a2.south) -- (b3.north);
\draw[<->, omCode] (a3.south) -- (b5.north);
\node[font=\scriptsize, align=left, anchor=west] at (6.6,-0.65) {8 steps: each advances\\one list, or both on a match};
\end{tikzpicture}
\par\smallskip{\small The AND of two words as a merge of their sorted lists of document numbers. The shaded documents
are in both. Document 8 is passed on the first list because the second has jumped from 4 to 12; documents 1, 4 and 15
are passed on the second. After 27 is compared with 30, the first list runs out and the merge stops.}
\end{omfigure}

\omcode{code/bachelor-1/30-project-text-search-engine/listings/query.c}{33}{56}{The AND query: look up every term, sort
the lists by length, and merge them into the running result}\label{lst:b1:project-text-search-engine:merge}

A word that occurs in no document has no \omterm{def:b1:project-text-search-engine:inverted-index}{posting list} at all; then no document can contain every word, and the
function stops at once (line 39). Which of the survivors are the best? A document that uses a query word often is
probably about it, and a rare word says more about a document than a word that every document uses. The usual score
combines the two.

\begin{definition}[Term frequency, document frequency, tf--idf score]\label{def:b1:project-text-search-engine:tfidf}
In a collection of $N$ documents, the \emph{term frequency}\index{term frequency} $\mathrm{tf}(t, d)$ of a word $t$ in a
document $d$ is the number of times $t$ occurs in $d$, and its \emph{document frequency}\index{document frequency}
$\mathrm{df}(t)$ is the number of documents that contain it. The \emph{tf--idf score}\index{tf--idf score} of a document
$d$ for a query of words $t_1, \dots, t_k$ is
\[
  \mathrm{score}(d) = \sum_{i=1}^{k} \mathrm{tf}(t_i, d) \cdot \ln \frac{N}{\mathrm{df}(t_i)} .
\]
\end{definition}

The factor $\ln (N / \mathrm{df})$ is the \emph{inverse \omterm{def:b1:project-text-search-engine:tfidf}{document frequency}}: zero for a word in every document, which
therefore cannot tell documents apart, and largest for a word in only one. It is a choice, not a theorem. Other
engines dampen the \omterm{def:b1:project-text-search-engine:tfidf}{term frequency} (a word used ten times is not ten times as relevant) or divide by the length of the
document, and their results differ; this one is the plain textbook form. Both factors are already in the index: the
\omterm{def:b1:project-text-search-engine:tfidf}{term frequency} is the count of a posting, and the \omterm{def:b1:project-text-search-engine:tfidf}{document frequency} is the length of the \omterm{def:b1:project-text-search-engine:inverted-index}{posting list}.

\begin{history}[Inverse document frequency]
% ledger: b1:tfidf
In 1972, Karen Spärck Jones published in the \emph{Journal of Documentation} ``A statistical interpretation of term
specificity and its application in retrieval''. She argued that a term matching many documents should count for
less than one matching few, and weighted each term by a function of the number of documents it occurs in. The
logarithmic weight became known as the inverse \omterm{def:b1:project-text-search-engine:tfidf}{document frequency}. Combined with the \omterm{def:b1:project-text-search-engine:tfidf}{term frequency}, it is still the
starting point of most ranking functions in search engines.
\end{history}

The query computes the score of each surviving document, finding each term's count in its sorted \omterm{def:b1:project-text-search-engine:inverted-index}{posting list} by
binary search, and keeps the best ten in a min-heap of size ten, the method of \cref{lst:b1:heaps:topk}: a new hit only
has to be compared with the worst hit kept, at the root.

\omcode{code/bachelor-1/30-project-text-search-engine/listings/topk.c}{34}{46}{Keeping the $k$ best hits: \omterm{def:b1:heaps:sift}{sift up} while
there is room, then replace the root only by a better hit}\label{lst:b1:project-text-search-engine:topk}

``Better'' is defined by \code{hit_better}: a higher score or, on equal scores, a smaller document number. Ties are
frequent here, since two documents in which the query words occur equally often get the same score, and without this
rule the order of equal hits would depend on the order of the heap's operations. The naive scan would then print the
same documents in another order, and the test that compares them would fail for a reason that is not an error.
With the rule, the order of the hits is a property of the hits alone.

\omoutput[26-41]{code/bachelor-1/30-project-text-search-engine/outputs/build.txt}{Two queries. Quixotic is in two
documents, once each: $\ln(400/2) = 5.2983$. Of the 34 documents with nodo and beli, the best uses nodo twice and beli
four times.}\label{out:b1:project-text-search-engine:queries}

The first score can be checked by hand: \emph{nodo} is in 45 documents and \emph{beli} in 265, so the best document
scores $2 \ln(400/45) + 4 \ln(400/265) = 4.3697 + 1.6469 = 6.0165$. The query \emph{Quixotic}, capital included, is cut
by the same \omterm{def:b1:project-text-search-engine:tokenizer}{tokenizer} and found.

\section{Tests, makefile and Git history}\label{sec:b1:project-text-search-engine:tests}

The project has five modules, three programs and one test program, so the \omterm{def:b1:c-modules-make-tests:makefile}{makefile} of
\cref{lst:b1:c-modules-make-tests:makefile} grows a little: one variable lists the core modules that every program
links, and each program adds what it needs.

\omcode{code/bachelor-1/30-project-text-search-engine/listings/Makefile}{1}{35}{The project's
\omterm{def:b1:c-modules-make-tests:makefile}{makefile}}\label{lst:b1:project-text-search-engine:makefile}

\omoutput[1-18]{code/bachelor-1/30-project-text-search-engine/outputs/build.txt}{A build from scratch, then the \omterm{def:b1:c-modules-make-tests:unit-test}{unit
tests}.}\label{out:b1:project-text-search-engine:make}

The \omterm{def:b1:c-modules-make-tests:unit-test}{unit tests} check each module on its own, on small inputs whose answers are known: the \omterm{def:b1:project-text-search-engine:tokenizer}{tokenizer} on the line of
\cref{lst:b1:project-text-search-engine:tok-test}; the index on four documents of a few words, then after 3000 more
words, which make the table grow; the heap on a stream of seven scores with a tie. They catch an error close to where
it is. They cannot show that the modules fit together, that the program reads the folder it is given, or that its
output is the one the specification asks for. That is the job of another kind of test.

\begin{definition}[End-to-end test]\label{def:b1:project-text-search-engine:end-to-end}
An \emph{end-to-end test}\index{end-to-end test} runs a whole program as its users do, from its real inputs (here a
folder and a command line) to its real outputs (the printed text and the \omterm{def:b1:the-shell:exit-status}{exit status}), and checks those outputs,
without looking inside the program.
\end{definition}

Here the \omterm{def:b1:project-text-search-engine:end-to-end}{end-to-end test} has an oracle: the naive scan, which shares the \omterm{def:b1:project-text-search-engine:tokenizer}{tokenizer}, the score and the printing with
the engine, but not the index. The test builds both programs with the \omterm{def:b1:c-modules-make-tests:makefile}{makefile} in a fresh temporary folder, generates
the corpus, and runs both on the 200 seeded queries of one to three words, plus the planted words, a query with a
capital and a punctuation mark, and a word found nowhere. For every query the two outputs and the two exit statuses
must be identical, byte for byte.

\omcode{code/bachelor-1/30-project-text-search-engine/tests/test_end_to_end.py}{15}{44}{The \omterm{def:b1:project-text-search-engine:end-to-end}{end-to-end test}: build,
generate, and compare with the oracle on every query}\label{lst:b1:project-text-search-engine:e2e}

The scores are printed with four decimals, so the comparison is exact only because both programs add the same
terms in the same order, through the same function \code{tfidf}; two different orders of addition can round
differently in the last bit. The test also builds the engine with the address and undefined-behaviour sanitizers and
requires them to stay silent on twenty queries, an empty query and a word found nowhere. That covers the paths that
free memory early, such as the word with no \omterm{def:b1:project-text-search-engine:inverted-index}{posting list}. Here is such a run by hand:

\omoutput{code/bachelor-1/30-project-text-search-engine/outputs/asan.txt}{The engine built with the sanitizers: the
same answers, and not a word on the error stream.}\label{out:b1:project-text-search-engine:asan}

\omoutput[42-49]{code/bachelor-1/30-project-text-search-engine/outputs/build.txt}{The two programs agree on a query, and
the exit statuses 1 (nothing found) and 2 (no such folder).}\label{out:b1:project-text-search-engine:status}

\paragraph{The history.} The project was kept in Git from its first milestone, with one \omterm{def:b1:version-control:commit}{commit} per milestone, so that each
\omterm{def:b1:version-control:commit}{commit} is a version that does one more thing. The session below replays that history in a fresh \omterm{def:b1:version-control:repository}{repository}: it copies
the files of each milestone and \omterm{def:b1:version-control:commit}{commits} them, then shows the log, the content of one \omterm{def:b1:version-control:commit}{commit}, and a \omterm{def:b1:version-control:staging}{working tree} with
nothing left to \omterm{def:b1:version-control:commit}{commit} after a build, since the \code{.gitignore} lists the programs, the \omterm{def:b1:compilation-pipeline:compiler}{object files} and the
generated corpus.

\omoutput{code/bachelor-1/30-project-text-search-engine/outputs/history.txt}{The project's history, one \omterm{def:b1:version-control:commit}{commit} per
milestone, and the clean \omterm{def:b1:version-control:staging}{working tree} after a build.}\label{out:b1:project-text-search-engine:history}

Generated files stay out of the \omterm{def:b1:version-control:repository}{repository}: the corpus is 1.9 megabytes that the script rewrites in a second, and a
committed copy could drift from the script that is supposed to make it.

\section{Measured against the naive scan}\label{sec:b1:project-text-search-engine:measure}

Is the index worth its cost? Building it reads every file once, as one naive query does, and also fills a \omterm{def:b1:hash-tables:hash-table}{hash table}.
After that, a query touches only the \omterm{def:b1:project-text-search-engine:inverted-index}{posting lists} of its words. A small program, \code{measure}, times the building of
the index, the 200 queries answered ten times over from the index, and the first twenty queries answered by the scan.
It checks that the two methods agree on those twenty, and prints the number of queries from which the index has paid
for itself.

\omoutput{code/bachelor-1/30-project-text-search-engine/outputs/timing.txt}{The index against the naive scan, measured on
the author's laptop.}\label{out:b1:project-text-search-engine:timing}

One query from the index takes about three microseconds, against nearly six milliseconds for the scan, which reads
1.9 megabytes and compares every word of them with every term: nearly two thousand times faster. The index costs about
two scans to build (more than one, since it also hashes every word and grows the table), so it pays for itself from the
third query on.

\begin{omfigure}
\begin{tikzpicture}
\begin{axis}[width=0.8\linewidth, height=5.6cm, xlabel={number of queries $q$}, ylabel={total time (ms)},
  xmin=0, xmax=10, ymin=0, ymax=60, legend pos=north west, legend style={font=\small, draw=none},
  tick label style={font=\small}, label style={font=\small}, grid=major, grid style={black!10}]
\addplot[omCode, very thick, mark=*, mark size=1.4pt] table[x=q, y=scan]
  {code/bachelor-1/30-project-text-search-engine/figures/measure.dat};
\addlegendentry{scan every file for every query}
\addplot[omThm, very thick, mark=square*, mark size=1.4pt] table[x=q, y=index]
  {code/bachelor-1/30-project-text-search-engine/figures/measure.dat};
\addlegendentry{build the index once, then query it}
\end{axis}
\end{tikzpicture}
\par\smallskip{\small Total time of $q$ queries on the corpus, from the measured output above. The scan's line rises by
almost six milliseconds per query; the index's line starts at its building time and stays flat, its queries far too
fast to show at this scale. The lines cross between the second and third queries.}
\end{omfigure}

The saving grows with the corpus. The scan's cost is proportional to the total size of the texts, while a query on the
index costs in proportion to the lengths of the \omterm{def:b1:project-text-search-engine:inverted-index}{posting lists} of its words. For the rare words that make good queries,
those lists are short, and Zipf's law says that most words are rare. The cost of the index is memory, a little more than
the texts here, and the time to build it again when a text changes. A real engine updates it piece by piece, which
this one cannot do.

\paragraph{What would come next.} Three improvements are classic. \emph{Stemming} would let \emph{index} match
\emph{indexes}, by cutting words to a common root, a rule that depends on the language. \emph{Phrase queries}, such as
the five words of a title in that order, need the positions of the words in the postings, and
\cref{exo:b1:project-text-search-engine:10} adds them. \emph{Compressing the postings} would make the index smaller
than the texts: the gaps between successive document numbers are small, and small numbers fit in a byte or less. The
algorithms on texts in the next volume go further, and \cref{exo:b1:project-text-search-engine:12} makes the index
survive the end of the program by writing it to a file.

\section{Exercises}\label{sec:b1:project-text-search-engine:exercises}

\begin{exercise}[$\star$]\label{exo:b1:project-text-search-engine:1}
Give the words that \code{next_word} (\cref{lst:b1:project-text-search-engine:next-word}) finds in the UTF-8 line
\texttt{It's 9:30-{}-{}Ünal's "e-{}mail" (v2.0)!}, in order. Which of them would a reader not call words, and why does
the query \emph{Ünal} still find this line?
\end{exercise}

\begin{exercise}[$\star$]\label{exo:b1:project-text-search-engine:2}
Three documents are indexed in this order: document 0, \emph{the cat sat on the mat}; document 1, \emph{The Cat and
the Hat}; document 2, \emph{a cat, a hat; a bat}. Write the \omterm{def:b1:project-text-search-engine:inverted-index}{posting lists} of \emph{the}, \emph{cat}, \emph{hat} and
\emph{a}, and give the number of distinct words and the number of postings of the index.
\end{exercise}

\begin{exercise}[$\star$]\label{exo:b1:project-text-search-engine:3}
Merge the \omterm{def:b1:project-text-search-engine:inverted-index}{posting lists} of documents $[2, 5, 9, 14, 20, 31]$ and $[1, 5, 6, 14, 31, 40, 41]$ as the loop of
\cref{lst:b1:project-text-search-engine:merge} does. Give the result and the number of times the loop's body runs.
Why does the loop stop before the end of the second list?
\end{exercise}

\begin{exercise}[$\star$]\label{exo:b1:project-text-search-engine:4}
A heap of size $k = 3$ (\cref{lst:b1:project-text-search-engine:topk}) receives hits of scores 4, 8, 1, 6, 9, 3, 7, for
documents 0 to 6 in this order. Give the heap's \omterm{def:b1:pointers-and-arrays:array}{array} after each push, and the three hits that \code{topk_sorted}
returns.
\end{exercise}

\begin{exercise}[$\star\star$]\label{exo:b1:project-text-search-engine:5}
A colleague counts words with the function below, a word being a run of ASCII letters. On the three texts of the
output it prints 4, 4 and 3. Explain the 3, give the class of input it fails on, and correct the function with one
line. Why does \code{next_word} not have this error?

\omcode{code/bachelor-1/30-project-text-search-engine/listings/exo_last_word.c}{5}{18}{Counting words at the
separators}
\omoutput[2-4]{code/bachelor-1/30-project-text-search-engine/outputs/exo_last_word.txt}{Three texts, the last without a
final newline.}
\end{exercise}

\begin{exercise}[$\star\star$]\label{exo:b1:project-text-search-engine:6}
The program below builds the \omterm{def:b1:project-text-search-engine:inverted-index}{posting list} of \emph{cat} for document 0, \emph{cat cat the}, and document 1, \emph{dog
cat}. Explain its output, including the negative inverse \omterm{def:b1:project-text-search-engine:tfidf}{document frequency}, and say which line of
\cref{lst:b1:project-text-search-engine:index} prevents the error in the engine. What would the error do to the
answers of the \omterm{def:b1:project-text-search-engine:end-to-end}{end-to-end test}?

\omcode{code/bachelor-1/30-project-text-search-engine/listings/exo_twice.c}{9}{20}{A \omterm{def:b1:project-text-search-engine:inverted-index}{posting list} that never looks at
its last posting}
\omoutput[2-3]{code/bachelor-1/30-project-text-search-engine/outputs/exo_twice.txt}{The \omterm{def:b1:project-text-search-engine:inverted-index}{posting list} of cat.}
\end{exercise}

\begin{exercise}[$\star\star$]\label{exo:b1:project-text-search-engine:7}
Write a function \code{query_or}, with the same arguments as \code{query_and}, that finds the documents containing
\emph{at least one} of the terms, scored by the sum of the tf--idf terms of the words they contain. Test it against a
brute-force count on small documents. How does its cost compare with that of \code{query_and}?
\end{exercise}

\begin{exercise}[$\star\star$]\label{exo:b1:project-text-search-engine:8}
Add an option \texttt{-{}-{}count} to \code{search}: \texttt{search DIR -{}-{}count WORD}\dots{} prints only the number of documents
that contain every word, with the same \omterm{def:b1:the-shell:exit-status}{exit status}. Write the counting function so that it allocates nothing and does
not compute a single score.
\end{exercise}

\begin{exercise}[$\star\star$]\label{exo:b1:project-text-search-engine:9}
An AND query has $k \ge 2$ words whose \omterm{def:b1:project-text-search-engine:inverted-index}{posting lists} have lengths $\ell_1 \le \ell_2 \le \dots \le \ell_k$.
\begin{enumerate}
\item Prove that the merges of \cref{lst:b1:project-text-search-engine:merge}, shortest list first, run their loop body
at most $(k - 1)\ell_1 + \ell_2 + \dots + \ell_k$ times in all.
\item Another method takes each document of the shortest list and looks for it in each other list by binary search.
Show that it costs $\Oh{\ell_1 \log \ell_k \cdot k}$ comparisons.
\item Compare the two for the query \emph{xylograph ba}, where \emph{ba} is in all 400 documents.
\end{enumerate}
\end{exercise}

\begin{exercise}[$\star\star\star$]\label{exo:b1:project-text-search-engine:10}
Add phrase queries: the documents in which the words $w_1, \dots, w_k$ occur consecutively, in this order. Store with
each occurrence its position in the document (its rank among the document's words), and find the positions $p$ such
that $w_i$ occurs at $p + i - 1$ for every $i$, by merging lists of (document, position) pairs. Test it on documents
where the words occur in both orders.
\end{exercise}

\begin{exercise}[$\star\star\star$]\label{exo:b1:project-text-search-engine:11}
Write the memory counted by \code{index_memory} as a function of the number of \omterm{def:b1:hash-tables:hash-table}{buckets} $B$, the number of distinct
words $W$, their total length $L$, the capacities $c_t$ of the \omterm{def:b1:project-text-search-engine:inverted-index}{posting lists}, and the documents' names, on a machine
where \omterm{def:b1:pointers-and-arrays:pointer}{pointers} take 8 bytes and \code{int} 4. On the corpus, $W = 5004$, $L = 30\,203$, $\sum_t c_t = 274\,086$, and the
400 names have 10 characters each. Find $B$ from the growth rule of the table, check the total of
\cref{out:b1:project-text-search-engine:stats}, and say what fraction of the posting \omterm{def:b1:pointers-and-arrays:array}{arrays} is slack.
\end{exercise}

\begin{exercise}[$\star\star\star$]\label{exo:b1:project-text-search-engine:12}
Make the index persist: write \code{index_save}, which writes the index into a \omterm{def:b1:c-strings-structs-files:text-binary}{binary file}, and \code{index_load},
which reads it back into an index that gives the same \omterm{def:b1:project-text-search-engine:inverted-index}{posting lists}, counts and document numbers. Specify the file's
format, make the loader refuse a truncated or corrupted file rather than crash, and write a round-trip test. Why is the
file not portable between machines as it stands?
\end{exercise}
